% Generated draft: requires mathematical and source review.
\documentclass[11pt]{article}
\usepackage[margin=25mm]{geometry}
\usepackage{fontspec}
\setmainfont{Latin Modern Roman}
\usepackage{amsmath,amssymb,mathtools,unicode-math}
\setmathfont{Latin Modern Math}
\usepackage{enumitem,needspace,xurl,hyperref}
\setlength{\parindent}{0pt}
\setlength{\parskip}{5pt}
\setlength{\emergencystretch}{3em}
\newcommand{\N}{\mathbb{N}}
\newcommand{\Z}{\mathbb{Z}}
\newcommand{\Q}{\mathbb{Q}}
\newcommand{\R}{\mathbb{R}}
\newcommand{\C}{\mathbb{C}}
\newcommand{\F}{\mathbb{F}}
\newcommand{\A}{\mathbb{A}}
\newcommand{\PP}{\mathbb{P}}
\newcommand{\E}{\mathbb{E}}
\newcommand{\eps}{\varepsilon}
\newcommand{\dd}{\,\mathrm{d}}
\begin{document}
\section*{UnsolvedMath ID: 701670}
\textbf{Source problem code:} catalogue entry
\subsection{Definitions and mathematical statement}
Let \(G\) be a group, with identity element \(1_G\), and let \(n\geq 1\) be an integer. Say that \(G\) has the \emph{unique \(n\)-fold product property}, denoted \(\mathrm{U}_n(G)\), if for every ordered tuple of finite nonempty subsets \(A_1,\ldots,A_n\subseteq G\), there exists \(g\in G\) such that
\[
\left|\left\{(a_1,\ldots,a_n)\in A_1\times\cdots\times A_n : a_1\cdots a_n=g\right\}\right|=1.
\]
The multiplication is taken in the displayed order. The usual unique product property is \(\mathrm{U}_2(G)\).

The source observes that \(\mathrm{U}_n(G)\) implies \(\mathrm{U}_m(G)\) whenever \(1\leq m\leq n\), by taking the remaining subsets to be \(\{1_G\}\).

\emph{Question (G. M. Bergman).} Is it true that, for every group \(G\) and every integer \(n\geq 2\),
\[
\mathrm{U}_2(G)\ \Longrightarrow\ \mathrm{U}_n(G)?
\]
Equivalently, are all the unique \(n\)-fold product conditions, for integers \(n\geq 2\), equivalent to the usual unique product condition?
\subsection{Short English statement}
If every pair of finite nonempty subsets of a group has a product element with exactly one representation, must the same hold for every ordered finite collection of at least two finite nonempty subsets?
\subsection{Sources}
\begin{itemize}
\item Local UnsolvedMath registry: permanent ID \(701670\), exact source problem code \texttt{catalogue entry}, title \emph{Kourovka Notebook Problem 20.12}. The supplied registry labels the problem open; that status is not independently verified here.
\item Supplied source attribution: \emph{The Kourovka Notebook (Unsolved Problems in Group Theory)}, 20th issue (2022), Problem 20.12; proposed by G. M. Bergman in 2022. The record also identifies the 21st issue edition (2026), \texttt{arXiv:1401.0300v48}.
\item The supplied collection description credits E. I. Khukhro and V. D. Mazurov as editors of Issues 1--21 (1965--2026).
\end{itemize}
\end{document}
